% !TEX program = xelatex
\documentclass[12pt,oneside]{report} % Thêm oneside để in 1 mặt

% --- CẤU HÌNH CỠ CHỮ VÀ LỀ TRANG ---
\usepackage[fontsize=13pt]{scrextend}
\usepackage[a4paper,top=40mm,bottom=25mm,left=40mm,right=25mm]{geometry}

% --- FONT CHỮ VÀ TIẾNG VIỆT/ANH ---
\usepackage{fontspec}
\setmainfont{Times New Roman}

% --- CẤU HÌNH CẤM NGẮT TỪ VÀ CHỐNG DÒNG MỒ CÔI (WIDOW/ORPHAN) ---
\hyphenpenalty=10000        % Cấm hoàn toàn ngắt từ thường
\exhyphenpenalty=10000     % Cấm ngắt từ cả ở các từ có sẵn dấu gạch nối
\binoppenalty=10000         % Cấm ngắt dòng tại các phép toán (+, -)
\relpenalty=10000           % Cấm ngắt dòng tại các dấu quan hệ (=, <, >)

% Chống tiêu đề/đoạn văn bị rớt đơn độc ở cuối trang
\clubpenalty=10000
\widowpenalty=10000
\displaywidowpenalty=10000

\usepackage{microtype}
\usepackage{setspace}

% --- CHÚ THÍCH (CAPTION) CHO BẢNG VÀ HÌNH ---
\usepackage{caption}
\captionsetup[table]{position=top, skip=6pt}    % Chú thích Bảng NẰM TRÊN
\captionsetup[figure]{position=bottom, skip=6pt} % Chú thích Hình NẰM DƯỚI

% --- TRÍCH DẪN VÀ TÀI LIỆU THAM KHẢO ---
\usepackage{biblatex}
\addbibresource{sections/ref.bib}

% --- CÁC GÓI BẢNG VÀ TRÌNH BÀY DỮ LIỆU ---
\usepackage{tabularx}
\usepackage{booktabs}
\usepackage{longtable}
\usepackage{array}

% --- CÁC GÓI ĐỒ HỌA VÀ HÌNH ẢNH ---
\usepackage{graphicx}
\usepackage{float}       % Quản lý vị trí hình/bảng [H]
\usepackage{subcaption}
\usepackage{eso-pic}
\usepackage{tikz}
\usetikzlibrary{positioning, shapes.geometric, arrows, arrows.meta}
\usepackage{pgfplots}
\pgfplotsset{compat=1.17}
\usepackage{lscape}

% --- TOÁN HỌC ---
\usepackage{amsmath}
\usepackage{amssymb}

% --- ĐỊNH DẠNG VĂN BẢN VÀ MÃ NGUỒN ---
\usepackage{indentfirst}
\usepackage{anyfontsize}
\usepackage{titlesec}
\usepackage{listings}
\usepackage{fancyhdr}
\usepackage{lastpage}

\lstset{
    basicstyle=\small\ttfamily,
    numbers=left,
    frame=single,
}

% --- CẤU HÌNH LIÊN KẾT (HYPERREF) ---
\usepackage{hyperref}
\hypersetup{
    colorlinks,
    citecolor=black,
    filecolor=black,
    linkcolor=black,
    urlcolor=black
}

% --- LỆNH TỰ ĐỊNH NGHĨA ---
\newcommand{\Khang}[1]{\textcolor{red!70!pink}{\textbf{\textit{#1}}}}

% --- PHÂN TRANG, THỤT ĐẦU DÒNG VÀ GIÃN DÒNG ---
\pagestyle{plain}             % Số trang ở cuối trang, canh giữa
\setlength{\parindent}{1.27cm} % Thụt đầu dòng 5 khoảng trắng (~1.27cm)
\setstretch{1.1}               % Giãn dòng 1.1 cho tài liệu Tiếng Anh

% ======================================================
% ĐỊNH DẠNG TIÊU ĐỀ (CHAPTER, SECTION, SUBSECTION)
% ======================================================

% Tiêu đề chương: Cỡ chữ 14, in đậm, chữ hoa, canh giữa
\titleformat{\chapter}[display]
  {\normalfont\fontsize{14pt}{16.8pt}\bfseries\centering}
  {\MakeUppercase{\chaptertitlename}\ \thechapter}
  {10pt}
  {\MakeUppercase}

% Tiêu đề Section: Cỡ chữ 13, in đậm, chữ thường
\titleformat{\section}
  {\normalfont\fontsize{13pt}{15.6pt}\bfseries}
  {\thesection}
  {1em}
  {}

% Tiêu đề Subsection: Cỡ chữ 13, in đậm, chữ thường
\titleformat{\subsection}
  {\normalfont\fontsize{13pt}{15.6pt}\bfseries}
  {\thesubsection}
  {1em}
  {}

\setlength{\emergencystretch}{2em}

% ======================================================
% BẮT ĐẦU TÀI LIỆU
% ======================================================
\begin{document}

\input{sections/cover-page}

\clearpage
\pagenumbering{roman} 
\input{sections/acknownledgement.tex}
\input{sections/abstract}

% --- THIẾT LẬP DÃN DÒNG 1.1 CHO TOÀN BỘ PHẦN MỤC LỤC VÀ CHƯƠNG ---
\setstretch{1.1}
\setlength{\parskip}{0pt}

% \input{sections/list-of-abbreviations}

\tableofcontents
\clearpage

\listoftables
\clearpage

\listoffigures
\clearpage

% ======================================================
% NỘI DUNG CHÍNH (Đánh số Ả Rập)
% ======================================================
\pagenumbering{arabic}
\setcounter{page}{1}

% -----------------------------------------
% CHƯƠNG 1
% -----------------------------------------
\chapter{Introduction}
\input{sections/1.Introduction/1.1.Motivation}
\input{sections/1.Introduction/1.2. Related_work}
\input{sections/1.Introduction/1.3.Scope}
\input{sections/1.Introduction/1.4.Thesis-structure}

% -----------------------------------------
% CHƯƠNG 2
% -----------------------------------------
\chapter{Theoretical background}
\label{chapter:theory}
\input{sections/2.Theoretical_background/2.1. Robot_Mecanum_Holonomic}
\input{sections/2.Theoretical_background/2.2.ArucoMarkerDetectionandPoseEstimation}
\input{sections/2.Theoretical_background/2.3.IMU_MPU6050}
\input{sections/2.Theoretical_background/2.4. Extended_Kalman_filter}
\input{sections/2.Theoretical_background/2.4.Matrix_Q_and_R}

% -----------------------------------------
% CHƯƠNG 3
% -----------------------------------------
\chapter{Hardware design}
\input{sections/3. Hardware design/3.1. Overall_System_Architecture}
\input{sections/3. Hardware design/3.2. Mecanum_Robot}
\input{sections/3. Hardware design/3.4. Esp32}
\input{sections/3. Hardware design/3.5.Marker_placement}

% -----------------------------------------
% CHƯƠNG 4
% -----------------------------------------
\chapter{Software design}
\input{sections/4. Software design/4.1. Communication_esp32_computer}
\input{sections/4. Software design/camera_calibration}
\input{sections/4. Software design/Aruco_marker_detection}
\input{sections/4. Software design/4.5. EKF_Sensor_Fusion – Detailed_Implementation}
\input{sections/4. Software design/Q_odometry}
\input{sections/4. Software design/R_IMU}
\input{sections/4. Software design/R_aruco_marker}

% -----------------------------------------
% CHƯƠNG 5
% -----------------------------------------
\chapter{Experiments and Evaluation}
\input{sections/5. Experiments and Evaluation/Experimental_setup}
\input{sections/5. Experiments and Evaluation/Result_analyse}
\input{sections/5. Experiments and Evaluation/summary}

% -----------------------------------------
% CHƯƠNG 6
% -----------------------------------------
\chapter{Conclusion and future work}
\input{sections/6. Conclusion and future work/6.1. CONCLUSION}
\input{sections/6. Conclusion and future work/6.2.Future_work}

% -----------------------------------------
% TÀI LIỆU THAM KHẢO
% -----------------------------------------
\input{sections/7. Reference/Reference}

% -----------------------------------------
% PHỤ LỤC
% -----------------------------------------
\chapter{Appendix}
\input{sections/8. appendices/appendix}

\end{document}